% !TITLE 木兰诗
% !AUTHOR 无名氏
% !DYNASTY 南北朝
% !GENRE 乐府民歌
% !ORDER 1

\begin{poem}
唧唧\textsuperscript{1}复唧唧，木兰当户\textsuperscript{2}织。不闻机杼\textsuperscript{3}声，唯闻女叹息。

问女何所思，问女何所忆。女亦无所思，女亦无所忆。昨夜见军帖\textsuperscript{4}，可汗\textsuperscript{5}大点兵，军书十二卷，卷卷有爷名。阿爷无大儿，木兰无长兄，愿为市鞍马\textsuperscript{6}，从此替爷征。

东市买骏马，西市买鞍鞯\textsuperscript{7}，南市买辔头\textsuperscript{8}，北市买长鞭。旦辞爷娘去，暮宿黄河边，不闻爷娘唤女声，但闻黄河流水鸣溅溅\textsuperscript{9}。旦辞黄河去，暮至黑山头，不闻爷娘唤女声，但闻燕山胡骑鸣啾啾\textsuperscript{10}。

万里赴戎机\textsuperscript{11}，关山度若飞。朔气传金柝\textsuperscript{12}，寒光照铁衣\textsuperscript{13}。将军百战死，壮士十年归。

归来见天子，天子坐明堂\textsuperscript{14}。策勋十二转\textsuperscript{15}，赏赐百千强\textsuperscript{16}。可汗问所欲，木兰不用尚书郎\textsuperscript{17}，愿驰千里足，送儿还故乡。

爷娘闻女来，出郭\textsuperscript{18}相扶将；阿姊闻妹来，当户理红妆\textsuperscript{19}；小弟闻姊来，磨刀霍霍\textsuperscript{20}向猪羊。开我东阁门，坐我西阁床，脱我战时袍，著我旧时裳。当窗理云鬓\textsuperscript{21}，对镜帖花黄\textsuperscript{22}。出门看火伴\textsuperscript{23}，火伴皆惊忙：同行十二年，不知木兰是女郎。

雄兔脚扑朔\textsuperscript{24}，雌兔眼迷离\textsuperscript{25}；双兔傍地走\textsuperscript{26}，安能辨我是雄雌？
\end{poem}

\begin{annotations}
\notetext{1}{唧唧：叹息声。也有人认为是织布机的声音。反复用“唧唧”，写出了木兰心中的忧愁。}
\notetext{2}{当户：对着门。木兰在门前织布，说明她是勤劳的农家女孩。}
\notetext{3}{机杼：织布机的梭子。不闻机杼声，意思是织布机没有响——木兰停下手中的活，在叹气。}
\notetext{4}{军帖：征兵文书。古代征兵的公文贴在城门上或送到各家。}
\notetext{5}{可汗：北方游牧民族对首领的称呼。这里指北魏的皇帝。北朝是鲜卑族政权，所以用“可汗”而非“天子”。}
\notetext{6}{市鞍马：购买马鞍和马匹。市，买。木兰要为父亲置办从军的装备。}
\notetext{7}{鞍鞯：马鞍和鞍垫。鞯（jiān），垫在马鞍下面的布垫。}
\notetext{8}{辔头：马笼头和缰绳。辔（pèi），用来控制马匹方向的器具。}
\notetext{9}{溅溅：流水声。黄河的水流声代替了爷娘的呼唤声，写出了征途的孤独。}
\notetext{10}{啾啾：战马的嘶鸣声。燕山胡人的马匹嘶叫，暗示已经接近战场。}
\notetext{11}{戎机：战争、军务。万里赴戎机，意思是奔赴万里之外的战场。}
\notetext{12}{金柝：古代军中打更用的铜器，又叫“刁斗”。朔气传金柝——北风传来了打更的声音，写出边塞夜晚的寒冷。}
\notetext{13}{铁衣：铁甲战袍。寒光照铁衣——寒冷的月光照在铁甲上，写出了边塞生活的艰苦。}
\notetext{14}{明堂：皇帝接见大臣、举行大典的殿堂。}
\notetext{15}{策勋十二转：记录功勋，连升十二级。转，军功每升一级叫一转。十二是虚数，形容功勋极高。}
\notetext{16}{百千强：形容赏赐极多。强，多的意思。}
\notetext{17}{尚书郎：尚书省的高级官员。木兰不稀罕当大官，只想回家，与世人追名逐利形成鲜明对比。}
\notetext{18}{出郭：走出城外。郭，外城。一家人远远地到城门外迎接，可见期盼之切。}
\notetext{19}{红妆：女子的装扮。阿姊听说妹妹回来，忙着梳妆打扮准备迎接。}
\notetext{20}{霍霍：磨刀的声音。小弟磨刀杀猪宰羊，合家欢庆的热闹场面。}
\notetext{21}{云鬓：像云一样蓬松柔美的鬓发。形容女子发式美丽。}
\notetext{22}{花黄：古代女子在额头上贴的黄色花饰。当窗理鬓、对镜贴黄——脱下战袍，木兰重新做回了女孩子。}
\notetext{23}{火伴：即“伙伴”，指同伍的士兵。古代军队十人为一“火”，同火吃饭，故称“火伴”。}
\notetext{24}{扑朔：扑腾、跳跃的样子。雄兔被抓住时会用力蹬腿挣扎。}
\notetext{25}{迷离：眯着眼睛、茫然的样子。雌兔被抓住时会眯缝起眼睛不动。}
\notetext{26}{傍地走：一起在地上跑。两只兔子一起奔跑时，你分不清哪只是雄哪只是雌——木兰以此幽默地反问：男女并肩做事，谁分得清男女呢？}
\end{annotations}

\begin{appreciation}
《木兰诗》是北朝民歌里最出色的叙事诗之一，也是中国文学史上最出名的女英雄故事。从迪士尼动画到真人电影，全世界都听过，最好的版本还是这首诗本身——它有电影拍不出来的东西。

开头是织布机边一声声叹息。家人问她想什么，她说没想什么——可昨夜军帖，卷卷都有我爹的名字。爹没有大儿子，我没有哥哥，那就买马买鞍，我替爹去。理由不是什么大口号，只是"阿爷无大儿，木兰无长兄"——眼看父亲要上战场，就自己顶上去。

"东市买骏马，西市买鞍鞯，南市买辔头，北市买长鞭"——她不会真跑四个市场，这是民歌的铺排，像说书人一口气念下去。最动人的细节在这："旦辞爷娘去，暮宿黄河边，不闻爷娘唤女声，但闻黄河流水鸣溅溅。"行军路上，听不见娘唤她的声音，只有黄河水哗哗响。后文写十二年战争只用了二十个字："将军百战死，壮士十年归。"刀光剑影一字没写，那十年有多苦，都写在那句"不闻爷娘唤女声"里了。

仗打完了，天子问她要什么，她不要官，只要快马回家。回家是全诗高潮——爹娘互相搀着出城接，姐姐梳妆，弟弟磨刀霍霍杀猪宰羊。木兰推开东阁门，坐上西阁床，脱掉战袍，穿上旧衣裳，对镜帖花黄——十二年过去，她还是那个爱美的女孩。

结尾拿兔子收场：雄兔脚扑朔，雌兔眼迷离，两只兔子一起跑，谁分得清公母？男女并肩做事，本不该分。

还记得《氓》吗？《氓》里的女子遇人不淑，最后"反是不思，亦已焉哉"；木兰是另一个答案——路自己选，日子自己过。《氓》在《卫风》，木兰诗在北朝，隔了一千多年，民歌讲故事的样子没变：不喊口号，只写日子。

妹妹，哥哥不指望你替谁从军。木兰有两件事值得学：家里有难处，你顶得上；认准的事，你走得稳。替哥哥挡挡妈妈的怒火，也要学着点。
\end{appreciation}
